\section{模型评价与推广}

\subsection{模型优势}

最终稿应结合实际完成情况评价以下方面：四问是否共享同一套状态和边界口径；变系数是否以通量形式保留；达标时刻是否由空间最大含水率首次穿越确定；问题四是否通过移动坐标与控制变量对照解释收缩影响。

\draftnote{只保留已由方程、代码和验证结果支持的优点。不要使用“精度很高”“适用性强”等没有证据的表述。}

\subsection{模型局限}

本题可能的主要局限包括：一维径向近似忽略端面和轴向差异；附件中的环境量未必等同于药材表面平衡含水率；经验物性公式的适用范围有限；半径曲线不能完整描述长度、孔隙率和组织结构变化。最终稿应说明哪些局限已经通过对照或敏感性分析评估，哪些仍构成结论边界。

\subsection{模型推广}

在获得其他药材的几何、平衡含水率、传热传质系数和物性参数后，本文的固定区域与移动区域框架可用于比较不同烘干条件下的内部均匀性和达标时长。若用于实际工艺优化，还需要加入能耗、有效成分保持率和实验标定数据，不能仅凭本题仿真直接给出生产参数。

\subsection{主要结论}

\begin{enumerate}
  \item 问题一：\resultblank{预热阶段最重要的温度与含水率结论}。
  \item 问题二：\resultblank{变物性全过程最重要的机制结论}。
  \item 问题三：固定半径条件下烘干时长为\resultblank{数值与单位}。
  \item 问题四：考虑收缩后的烘干时长为\resultblank{数值与单位}，收缩和物性变化的影响分别为\resultblank{对照结论}。
\end{enumerate}

